% BEGIN GENERATED ARTIFACT: prf:dnf-representation-of-boolean-functions-propositional-logic
\newpage
\phantomsection
\label{prf:dnf-representation-of-boolean-functions-propositional-logic}
\LRAProofFor{thm:dnf-representation-of-boolean-functions-propositional-logic}

\begin{remark*}[Return]
\hyperref[thm:dnf-representation-of-boolean-functions-propositional-logic]{Return to Theorem (DNF Representation of Boolean Functions).}
\end{remark*}

\begin{theorem*}[DNF Representation of Boolean Functions]
Every Boolean function \(f:\mathbb{B}^n\to\mathbb{B}\) is represented by a
formula in disjunctive normal form using variables \(P_1,\ldots,P_n\).
\end{theorem*}

\begin{proof}
\textbf{Professional Standard Proof.}~
TODO: supply the compact proof.
\end{proof}

\begin{proof}
\textbf{Detailed Learning Proof.}~
TODO: supply the detailed learning proof.
\end{proof}

\begin{remark*}[Proof structure]
Use the canonical disjunctive normal form theorem to construct, from the truth
table of \(f\), a disjunction of row conjunctions for exactly those rows on
which \(f\) has value \(\mathrm{True}\). The representation definition then
identifies the resulting DNF formula with the Boolean function on every input
row.
\end{remark*}

\begin{dependencies}
\begin{itemize}
  \item \hyperref[def:boolean-function-representation-propositional-logic]{Representation of a Boolean Function}
  \item \hyperref[thm:canonical-disjunctive-normal-form-propositional-logic]{Canonical Disjunctive Normal Form Theorem}
  \item \hyperref[def:disjunctive-normal-form-propositional-logic]{Disjunctive Normal Form}
\end{itemize}
\end{dependencies}

\clearpage
% END GENERATED ARTIFACT: prf:dnf-representation-of-boolean-functions-propositional-logic
